\subsection{Exponential Models and Change in Time} 
Another natural way to describe an exponential model is to choose a nice multiplier, and describe the amount of time it takes. We may also be given a value other than the starting value.
%\begin{example}
%A scientist culturing bacteria in a petri dish observes that the population doubles every 2.3 hours. If there are initially 300 bacteria in the dish, write a function describing the population $p$ of bacteria in the dish after $t$ hours.
	%\begin{lpic}{}{5}
		%\regtext{3}{-1}{Time (hours)};
		%\regtext{7}{-1}{Population};
		%\regtext{12}{-1}{\# of times doubled};
		%\foreach \x in {5,9} \draw (\x,0) -- (\x,-5);
		%\draw (1,-1) -- (15,-1);
	%\end{lpic}
	%\[
	%p(t)=\makeblank
	%\]
%\end{example}
%We may be given some other value instead of the starting value.
\begin{example}
A scientist culturing bacteria in a petri dish observes that the population doubles every 2.3 hours. After 10 hours, there are 10,000 bacteria in the dish. Write a function describing the population $p$ of bacteria in the dish after $t$ hours.
	\begin{lpic}{}{6}
		\regtext{3}{-1}{Time (hours)};
		\regtext{7}{-1}{Population};
		\regtext{12}{-1}{\# of times doubled};
		\foreach \x in {5,9} \draw (\x,0) -- (\x,-6);
		\draw (1,-1) -- (15,-1);
	\end{lpic}
	\[
	p(t)=\makeblank
	\]
\end{example}
\begin{displaybox}[Exponential Models (General Form)]
An exponential model is a function in the form
\[
f(x)=f_1\cdot m^{\frac{x-x_1}{\Delta x}}
\]
where
\begin{itemize}
\item $f_1$ is the value of $f$ when $x=x_1$.
\item $\Delta x$ (``delta $x$'') is the change in $x$ which results in the value of $f$ being multiplied by $m$.
\end{itemize}
\end{displaybox}
The multiplier may not be explicitly given, but this is easy to find.
\begin{example}
A scientist culturing bacteria in a petri dish observes that, after 5 hours, there are 3000 bacteria in the dish, and after 12 hours, there are 30,000 bacteria. Write a function describing the population $p$ of bacteria in the dish after $t$ hours.
	\begin{lpic}{}{4}
		\regtext{3}{-1}{Time (hours)};
		\regtext{7}{-1}{Population};
		\draw (5,0) -- (5,-3);
		\draw (1,-1) -- (9,-1);
	\end{lpic}
	\[
	p(t)=\makeblank
	\]
\end{example}
